\chapter[Hyperbolic components beyond the quadratic family]{Source digests VI: hyperbolic components beyond the quadratic family}
\label{ch:digests-sierpinski}

This chapter digests the paper of Dudko and Luo~\cite{DudkoLuoSierpinski} on hyperbolic
components of rational maps whose Julia sets are Sierpi\'nski carpets.  It is the
nearest analogue, in a family of rational maps, of the question that the survey~\cite{DudkoMLC}
lists as open for the Mandelbrot set: uniform bounded geometry of hyperbolic components.
Statements are those of the source, up to notation; numbering is that of the source.

\section{Definitions}
\label{sec:digest-sierpinski-defs}

\begin{definition}[Hyperbolic maps and components]
\label{def:sierp-hyperbolic}
A rational map $f$ of degree $d$ is hyperbolic if every critical point converges to an attracting periodic cycle.
In the fixed-point-marked moduli space $\mathcal M^{\rm fm}_d=\operatorname{Rat}^{\rm fm}_d/\operatorname{PSL}_2(\C)$, a connected component
of the open set of hyperbolic maps is a hyperbolic component $H$.
\end{definition}

\begin{definition}[Disjoint type]
\label{def:sierp-disjoint}
$H$ is of disjoint type if for every $[f]\in H$ each grand orbit of a Fatou component contains a unique critical orbit.
Equivalently, every $f\in H$ has exactly $2d-2$ attracting periodic cycles.
\end{definition}

\begin{definition}[Sierpi\'nski carpet component]
\label{def:sierp-carpet}
$H$ is a Sierpi\'nski carpet component if $J_f$ is homeomorphic to a Sierpi\'nski carpet for every $[f]\in H$.
\end{definition}

\begin{definition}[Multiplier map]
\label{def:sierp-multiplier}
For $H$ of disjoint type the map $\rho\colon H\to\mathbb D^{\,2d-2}$ records the multipliers of the $2d-2$ attracting cycles.
\end{definition}

\begin{definition}[Degeneration bounded by $K$]
\label{def:sierp-bounded-degeneration}
$[f]$ has degeneration bounded by $K$ if it has $K$-quasiconformal pseudo-Siegel disks $\widehat Z_{i,j}$ and
$K$-quasiconformal valuable-attracting domains $\widehat D_{i,j}$, and its pseudo-core surface
$\widehat X_f=\widehat\C\setminus\bigl(\bigcup\operatorname{Int}\widehat D_{i,j}\cup\bigcup\operatorname{Int}\widehat Z_{i,j}\bigr)$
satisfies $W_{\rm arc}(\widehat X_f)\le K$ and $W_{\rm loop}(\widehat X_f)\le K$.
\end{definition}

\begin{definition}[Regular and exceptional boundary; obstructed boundary]
\label{def:sierp-boundary}
$\partial H=\partial_{\rm reg}H\sqcup\partial_{\rm exc}H$, where $[f]\in\partial_{\rm reg}H$ if $f$ has exactly $2d-2$
non-repelling cycles, and $[f]\in\partial_{\rm exc}H$ otherwise, when at least two of these cycles collide.  The
obstructed boundary is $\partial^\infty H=\{\rho\in\partial\mathbb D^{2d-2}:\ \exists[f_n]\in H,\ [f_n]\to\infty,\ \rho([f_n])\to\rho\}$.
\end{definition}

\begin{remark}
The paper states the boundedness question as Conjecture~1.1, that every Sierpi\'nski carpet hyperbolic
component is bounded in $\mathcal M^{\rm fm}_d$.  It uses ``bounded'' in the ambient moduli space and gives no further
formal definition.
\end{remark}

\section{Results}
\label{sec:digest-sierpinski-results}

\begin{theorem}[Dudko--Luo, Theorem~A; parameter space]
\label{thm:sierp-A}
If $H\subset\mathcal M^{\rm fm}_d$ is a Sierpi\'nski carpet hyperbolic component of disjoint type, then $H$ is bounded in
$\mathcal M^{\rm fm}_d$, and $\rho\colon H\to\mathbb D^{2d-2}$ is an isomorphism that extends naturally.  In particular
$\partial H$ is locally connected.  Status: proved in the paper.
\end{theorem}

\begin{theorem}[Dudko--Luo, Theorem~B; dynamical plane]
\label{thm:sierp-B}
If $H$ is of Sierpi\'nski carpet and disjoint type, there is $\varepsilon>0$ such that every $[f]\in H$ and every non-repelling
periodic point $x$ of period $p$ admit a quadratic-like restriction $f^p\colon U\to V$ with $x\in U\subset V$ and
$\modulus(V\setminus U)\ge\varepsilon$.  Status: proved in the paper.
\end{theorem}

\begin{theorem}[Dudko--Luo, Theorem~C; a priori bounds]
\label{thm:sierp-C}
If $H$ is of disjoint type and $[f]\in\partial H$ is eventually golden mean, then $[f]$ has degeneration bounded by a
constant $K$ depending on $H$, on the pulled-off constant $N_{\rm Siegel}([f])$, and on the attracting multipliers.  Status: proved.
\end{theorem}

\begin{theorem}[Dudko--Luo, Theorem~1.4; precompactness]
\label{thm:sierp-precompact}
If $[f_n]\in\partial H$ are eventually golden mean with degeneration bounded by $K$, then a subsequence converges to some
$[f]\in\mathcal M^{\rm fm}_d$, and $[f]$ has $2d-2$ non-repelling cycles.  Status: proved.
\end{theorem}

\begin{theorem}[Dudko--Luo, Theorem~D; bi-neutral quadratics]
\label{thm:sierp-D}
For $f\in\partial_{\rm reg}H_{z^2}$ with two neutral fixed points and critical postcritical sets $P_1,P_2$, the maps
$f|_{P_1}$ and $f|_{P_2}$ have degree $1$, and $P_1,P_2$ are separated by an annulus whose modulus depends only on the
distance between the rotation multipliers.  Status: proved.
\end{theorem}

\begin{proposition}[Dudko--Luo, Proposition~2.2; multiplier embedding]
\label{prop:sierp-2-2}
$\rho\colon\partial_{\rm reg}H\to\partial\mathbb D^{2d-2}$ is an embedding, and $\rho(\partial_{\rm reg}H)\cap\rho(\partial_{\rm exc}H)=\varnothing$.
Status: proved, using multiplier transversality from the literature cited there.
\end{proposition}

\section{Relation to the Mandelbrot set and MLC}
\label{sec:digest-sierpinski-mlc}

The paper states that suitable a priori bounds for all quadratic polynomials are expected to imply MLC, citing
\cite{DudkoLyubichNeutral} for this heuristic, and it records progress for the primitive case through earlier work of
Kahn and Lyubich.  For the rational-map result, quadratic-like restrictions and the Douady--Hubbard straightening
theorem give a homeomorphism $R_{\rm ql}\colon H\to\Delta\subset\operatorname{Mand}^{\,2d-2}$, where $\Delta$ is the main
hyperbolic component of the Mandelbrot set.  The paper presents this as a dynamical refinement of
Theorem~\ref{thm:sierp-A}, not as a proof of MLC.

\section{Dependencies}
\label{sec:digest-sierpinski-deps}

The argument uses the pseudo-Siegel and renormalization bounds of the Dudko--Lyubich and Dudko--Luo--Lyubich papers;
the quasi-additivity law and covering lemma of Kahn--Lyubich; the density results cited for geometrically finite
degenerations; the postcritically finite Sierpi\'nski model of Pilgrim; expanding models of Dudko--Lyubich--Selinger;
multiplier transversality from the literature cited in the paper; and the Douady--Hubbard straightening theorem.
